\section{결합 규칙과 파라미터의 추가 분석}
\label{sec:appendix:fusion}

본문의 분석을 보완하기 위해 계층 가중치와 해상도 집계 대안, 보정과 집계의 메커니즘 분석을 제시한다.
계층 가중치와 집계 대안 표는 MVTec AD 2 \texttt{test\_public}, seed 0의 8개 클래스 macro이며, 메커니즘 분석은 별도로 명시한 3-seed 결과를 사용한다.

\subsection{계층 가중치}
최종 구성(Q95, Rényi)에서 $z_d=(z_{11}+z_{23})/2$에 더하는 얕은 계층 가중치 $\lambda$를 바꾸어 비교한다(\cref{tab:lambda_final}).
$\lambda=0.5$에서는 세 계층의 계수가 같다. $\lambda=0$은 얕은 계층의 점수 기여만 제외한다.

\begin{table}[!htbp]
\centering\small
\caption{얕은 계층 가중치 $\lambda$의 민감도 (MVTec AD 2, seed 0, Q95, Rényi, \AUPRO, \%). 굵은 값은 최고값이다.}
\label{tab:lambda_final}
\begin{tabular*}{\columnwidth}{@{\extracolsep{\fill}}lrr}
\toprule
$\lambda$ & 2-shot & 16-shot \\
\midrule
0 & 59.05 & 59.11 \\
0.0625 & 59.76 & 59.73 \\
0.125 (사용값) & \textbf{59.92} & 60.16 \\
0.25 & 59.32 & \textbf{60.36} \\
0.5 & 56.23 & 59.47 \\
\bottomrule
\end{tabular*}
\end{table}

얕은 계층을 제외하거나($\lambda=0$) 가중치를 0.5로 높이면 두 shot 모두 성능이 낮아진다. 16-shot에서는 0.25가 0.125보다 0.20\,pp 높아, 최적값은 shot에 따라 조금 다르다.

\subsection{해상도 집계 대안}
\label{sec:appendix:aggregation}
최종 구성의 특징, 뱅크, Q95 보정, 계층 가중치를 고정하고 해상도 집계 규칙만 바꾸었다(\cref{tab:aggregation_alternatives}).
선형 순위 가중치는 $(0.1, 0.2, 0.3, 0.4)$다.
두 순위 가중 집계는 Mean보다 \AUPRO가 높고, Rényi와 선형 가중의 차이는 0.1\,pp 이내다.
따라서 이 결과는 높은 순위에 큰 가중치를 주는 집계의 효과를 보여주며, 특정 Rényi 가중치가 다른 순위 가중치보다 뚜렷하게 낫다는 근거는 아니다.

\begin{table}[!htbp]
\centering\scriptsize\setlength{\tabcolsep}{3pt}
\caption{해상도 집계 대안 (MVTec AD 2, seed 0, Q95, \%). 각 셀은 \AUPRO{} / \pAUROC{}이며, 굵은 값은 열별 최고값이다.}
\label{tab:aggregation_alternatives}
\begin{tabular*}{\columnwidth}{@{\extracolsep{\fill}}lcc}
\toprule
집계 방식 & 2-shot & 16-shot \\
\midrule
Mean & 59.33 / \textbf{76.48} & 59.57 / 77.46 \\
Order (linear) & 59.91 / 76.45 & 60.08 / \textbf{77.55} \\
Order (Rényi, Ours) & \textbf{59.92} / 76.36 & \textbf{60.16} / 77.52 \\
\bottomrule
\end{tabular*}
\end{table}

\subsection{보정과 Rényi 집계의 메커니즘 분석}
\label{sec:appendix:mechanism}

\begin{figure*}[t]
\centering
\includegraphics[width=\textwidth]{figures/fig_scopead_calibration_mechanism.pdf}
\caption{\textbf{Q95 보정과 Rényi 집계의 결합 메커니즘 분석.}
MVTec AD 2, 2/16-shot, seed 0/1/2에서 같은 특징과 거리를 재사용했다.
(a) 상대 추가 응답 $E/\mu$의 Q95$-$None 변화. Normal은 정상 이미지, Background는 결함 이미지의 정상 배경, Defect는 결함 영역이다. 같은 픽셀 표본으로 조건 간 차이를 계산하고 이미지·seed·클래스 순으로 평균했다. 작은 점은 3-seed 평균 클래스 효과, 마름모와 오차 막대는 macro와 클래스 bootstrap 95\,\% 구간이다.
(b) None+Rényi 대비 순위만, 점수 크기만, 둘 다 보정한 경우의 \AUPRO{} 차이. 작은 점과 구간은 (a)와 같다.
(c) Q95 분모에 평균 1의 log-normal 오차를 넣었을 때 오차가 없는 경우 대비 \AUPRO{} 변화. 각 $\sigma$에서 seed별 3회 오차를 두 집계에 동일하게 주고 평균했다. 색은 shot, 점선/실선은 Mean/Rényi다.}
\label{fig:calibration-mechanism}
\end{figure*}

\noindent\textbf{상대 추가 응답($E/\mu$).}
Rényi 집계에서 정상 이미지의 $E/\mu$는 Q95 보정 후 2-shot에서 $7.74\!\rightarrow\!7.10$\,\%, 16-shot에서 $8.56\!\rightarrow\!7.55$\,\%로 감소했다(\cref{fig:calibration-mechanism}(a)).
결함 영역에서는 $5.34\!\rightarrow\!5.36$\,\%, $6.36\!\rightarrow\!6.11$\,\%로 변화가 작았다.
결함과 정상 배경의 $E/\mu$ 차이는 $+0.74/+0.80$\,pp 커졌고, 정상 이미지 $E/\mu$의 Q99는 $-1.30/-1.61$\,pp 줄었다.
이는 상대 추가 응답 분포의 변화이며, 고정 임계값에서의 오경보율 변화를 직접 측정한 것은 아니다.

\noindent\textbf{순위 변경과 점수 크기 변경의 분리.}
Q95로 순위만 정하고 원점수를 합산하는 조건과, 원점수 순위를 유지한 채 Q95 점수를 합산하는 조건을 비교했다(\cref{fig:calibration-mechanism}(b)).
순위만 바꾼 효과는 $+0.05/-0.14$\,pp인 반면, 크기만 바꾼 효과는 $+0.46/+0.70$\,pp로 최종 이득의 대부분을 차지했다.
따라서 보정의 효과는 주로 점수 크기의 조정에서 나온다.

\noindent\textbf{보정 오차에 대한 민감도.}
보정 분모에 평균 1의 log-normal 오차를 넣으면(\cref{fig:calibration-mechanism}(c)), $\sigma=0.3$에서 Mean의 \AUPRO 변화는 $-0.03/+0.06$\,pp인 반면 Rényi는 $-0.74/-0.47$\,pp였다.
Rényi 집계는 Mean보다 분모의 정확도에 더 민감하다.

\noindent\textbf{단위에 따른 분포 비교.}
$E/\mu$는 해당 픽셀의 평균 응답을 분모로 쓰므로 결함에서 $\mu$가 커지는 효과도 포함한다.
이를 확인하기 위해 다른 fold 정상 이미지 $\mu$의 Q95로 나눈 $E$도 측정했다.
이 값의 정상/결함 평균은 2-shot에서 $4.31\!\rightarrow\!3.80$ / $6.18\!\rightarrow\!6.19$\,\%, 16-shot에서 $4.57\!\rightarrow\!3.80$ / $8.01\!\rightarrow\!7.66$\,\%로, 정상 추가 응답이 줄고 결함 응답이 상대적으로 유지되는 방향은 같았다.
결함$-$배경 $E/\mu$ 차이 변화의 클래스 bootstrap 95\,\% 구간은 $[+0.06,+1.72]/[+0.00,+2.02]$\,pp이며, 16-shot 하한이 0에 가깝다.
정상 이미지 $E/\mu$의 Q99 변화 구간은 $[-2.29,-0.50]/[-3.00,-0.48]$\,pp이고, 감소한 클래스는 7/8과 6/8이다.

\begin{table*}[t]
\centering\small\setlength{\tabcolsep}{8pt}
\caption{\textbf{클래스별 검출 변화와 해상도 응답 집중도의 연관.} MVTec AD 2, 3-seed 평균. Coverage 상호작용은 $(\mathrm{Q95{+}R}-\mathrm{Q95{+}Mean})-(\mathrm{None{+}R}-\mathrm{None{+}Mean})$이며, GT 영역 coverage를 이미지 안에서 평균한 뒤 이미지와 seed를 평균한 값이다. 임계값은 정상 이미지의 교차검증으로 목표 FPR 1\,\%에 맞추었으며, 실제 평가 FPR은 조건마다 다르다. $\rho$는 Q95 상태의 응답 집중도와 Mean$\rightarrow$Rényi coverage 차이의 클래스 내 Spearman 상관이다.}
\label{tab:mechanism_class_effects}
\begin{tabular}{l r r r r}
\toprule
클래스 & 2-shot 상호작용 (pp) & 16-shot 상호작용 (pp) & 2-shot $\rho$ & 16-shot $\rho$ \\
\midrule
can & $-0.05$ & $-0.06$ & 0.48 & 0.39 \\
fabric & $+7.24$ & $+11.50$ & 0.56 & 0.31 \\
fruit\_jelly & $+0.23$ & $+0.72$ & 0.48 & 0.46 \\
rice & $-0.38$ & $-0.71$ & 0.38 & 0.22 \\
sheet\_metal & $-1.86$ & $-1.78$ & 0.48 & 0.43 \\
vial & $+0.12$ & $-0.18$ & 0.30 & 0.03 \\
wallplugs & $-0.19$ & $-0.59$ & 0.51 & 0.52 \\
walnuts & $+0.05$ & $+0.79$ & 0.29 & 0.30 \\
\midrule
Macro & $+0.65$ & $+1.21$ & --- & --- \\
fabric 제외 & $-0.30$ & $-0.26$ & --- & --- \\
\bottomrule
\end{tabular}
\end{table*}

\noindent\textbf{클래스와 이미지의 불확실성.}
\cref{tab:mechanism_class_effects}의 coverage 상호작용은 클래스에 따라 부호가 다르다.
클래스 bootstrap 구간은 0을 포함하지만, 서포트·임계값·클래스를 고정하고 이미지만 재표집한 구간은 $[+0.18,+1.16]/[+0.61,+1.92]$\,pp다.
두 구간은 서로 다른 변동을 나타내며, 후자는 클래스 구성이나 서포트 선택의 변동을 포함하지 않는다.

\noindent\textbf{결함 크기와 대비.}
GT 면적이 이미지의 0.1\,\% 미만인 작은 영역의 coverage 상호작용은 $+0.86/+1.53$\,pp지만 fabric을 제외하면 $-0.18/-0.06$\,pp다.
0.1--1\,\% 및 1\,\% 이상 구간은 can이 없어 7개 클래스 범위다.
회색조 대비(GT 영역 평균과 8픽셀 주변 정상 영역의 차이를 주변 표준편차로 나눈 값)와 결함 크기의 상관 방향은 클래스마다 달랐다.
따라서 작은 결함이나 저대비 결함 전반에서 이득이 있다고 보기는 어렵다.

\noindent\textbf{해상도 순위 점유율.}
\cref{fig:order-rank-diagnostic}는 seed 0에서 각 픽셀의 최고 점수를 낸 해상도의 비율이다.
Q95 보정 후 448의 정상 점유율은 47.5/50.7\,\%에서 33.9/36.4\,\%로 줄고, 16-shot 작은 결함의 1120 점유율은 14.5\,\%에서 29.2\,\%로 늘었다.
정상 픽셀의 1120 점유율도 13.7\,\%에서 22.3\,\%로 늘었으므로, 이 그림은 순위 편중의 변화를 보여주는 보조 자료로 사용한다.

\begin{figure}[tbp]
\centering
\includegraphics[width=\columnwidth]{figures/fig_scopead_order_rank_diagnostic.pdf}
\caption{\textbf{정상 참조 보정과 해상도 순위 점유율.} MVTec AD 2, seed 0, 8개 클래스 macro. 점선은 균등 점유율 25\,\%다.}
\label{fig:order-rank-diagnostic}
\end{figure}

\subsection{계층·해상도 보정 축의 분리}
\label{sec:appendix:calibration-axis}

\begin{table}[!htbp]
\centering\scriptsize\setlength{\tabcolsep}{3pt}
\caption{\textbf{Q95 보정 분모의 계층·해상도 성분 분리.} MVTec AD 2, 2/16-shot, seed 0/1/2 평균 \AUPRO{} (\%). Global은 배그별 공통 배율만, Layer/Resolution은 공통 배율에 해당 축 성분만, Separable은 두 축 성분을 함께 적용한다. Full은 원래 Q95다.}
\label{tab:calibration_axis}
\begin{tabular*}{\columnwidth}{@{\extracolsep{\fill}}lrrrr}
\toprule
& \multicolumn{2}{c}{2-shot} & \multicolumn{2}{c}{16-shot} \\
\cmidrule(lr){2-3}\cmidrule(lr){4-5}
분모 & Mean & Rényi & Mean & Rényi \\
\midrule
None & 56.12 & 56.32 & 59.73 & 59.72 \\
Global & 56.12 & 56.32 & 59.73 & 59.72 \\
Layer & 55.91 & 56.13 & 59.90 & 59.82 \\
Resolution & 56.26 & 56.87 & 59.89 & 60.55 \\
Separable & 56.02 & 56.78 & 60.02 & 60.73 \\
Full Q95 & 56.03 & 56.78 & 60.06 & 60.74 \\
\bottomrule
\end{tabular*}
\end{table}

\noindent\textbf{분모 분해.}
정상 서포트 Q95 $c_{l,r,b}$의 로그를 각 배그에서 다음과 같이 분해한다:
\begin{equation}
  \log c_{l,r,b}=\gamma_b+u_{l,b}+v_{r,b}+h_{l,r,b}.
  \label{eq:calibration-axis-factor}
\end{equation}
$\gamma_b$는 계층·해상도 전체 평균이고, $u_{l,b}$와 $v_{r,b}$는 각 축 평균에서 $\gamma_b$를 뺀 값이며, $h_{l,r,b}$는 잔여 항이다.
Global, Layer, Resolution, Separable은 각각 $\exp(\gamma_b)$, $\exp(\gamma_b+u_{l,b})$, $\exp(\gamma_b+v_{r,b})$, $\exp(\gamma_b+u_{l,b}+v_{r,b})$를 분모로 쓴다. Full은 원래 $c_{l,r,b}$이며 None은 1이다.
모든 축 개입은 배그 공통 척도를 유지하므로 Global을 대조로 사용한다.
특징, 뱅크, 계층 가중치, FP16을 고정하고 분모만 바꾸었으며(8개 클래스 $\times$ 2/16-shot $\times$ 3 seed), None과 Full의 결과는 \cref{tab:interaction_multiseed}와 일치한다.

\noindent\textbf{축 효과.}
Rényi 아래에서 Resolution$-$Global은 $+0.55/+0.82$\,pp이며, seed별 macro 효과는 2-shot $+0.56/+0.64/+0.45$\,pp, 16-shot $+0.82/+0.78/+0.87$\,pp로 모두 양수다.
Mean 아래의 같은 효과는 $+0.15/+0.15$\,pp로 작다.
Layer$-$Global은 Rényi 아래에서 $-0.19/+0.10$\,pp로 shot에 따라 방향이 달랐고, Full$-$Separable은 두 shot 모두 0.01\,pp 이내로 잔여 항의 macro 영향은 작았다.
다만 Resolution$-$Global에서 양의 효과를 보인 클래스는 두 shot 모두 8개 중 4개이고, fabric의 효과가 $+5.18/+7.26$\,pp로 크다.
클래스 bootstrap 95\,\% 구간은 $[-0.36,+1.99]/[-0.40,+2.80]$\,pp로 0을 포함하며, fabric을 제외하면 $-0.11/-0.10$\,pp다.

\subsection{평균 응답 조건부 분석과 영역·픽셀 가중}
\label{sec:appendix:conditional-mechanism}

\noindent\textbf{같은 평균 응답 구간에서의 비교.}
두 보정 조건의 같은 픽셀을 보정 전 평균 응답 $\mu$의 10분위 구간에 배정했다. 구간 경계는 다른 정상 fold에서 정했다.
이미지$\rightarrow$seed$\rightarrow$클래스 순으로 집계하면 정상과 배경의 Q95$-$None $E/\mu$ 변화는 두 shot 모두 10개 구간 전체에서 음수였다.
결함$-$배경 변화는 8개 클래스가 모두 관측된 구간(2-shot 5--9, 16-shot 6--9)에서 모두 양수였다.
따라서 평균 응답 수준이 다른 픽셀이 섞여서 생긴 현상으로 보기는 어렵다. 다만 보정 후의 $\mu$까지 고정한 비교는 아니다.

\noindent\textbf{Fabric의 영역 수와 결함 픽셀 수.}
Fabric의 GT 영역 192개 중 180개(93.75\,\%)는 이미지 면적의 0.1\,\% 미만이지만, 이 영역들은 전체 결함 픽셀의 1.39\,\%만 차지한다.
목표 FPR 1\,\% 임계값에서 Q95+Rényi$-$None+Rényi의 영역 평균 coverage 변화는 작은 영역에서 $+6.85/+10.41$\,pp, 나머지 영역에서 $-1.36/-0.84$\,pp다(2/16-shot).
같은 변화를 GT 픽셀 수로 가중하면 작은 영역의 TPR 기여는 $+0.06/+0.08$\,pp, 나머지는 $-2.66/-1.64$\,pp다.
즉 많은 작은 영역의 coverage 증가와 큰 영역의 손실이 함께 나타나며, 영역 단위 지표(\AUPRO)와 픽셀 단위 지표(\pAUROC)가 다르게 움직이는 이유를 설명한다.
